\clearpage\pagestyle{otMain}\setlength{\otBlocEspacement}{2.5mm}\fontsize{11}{15}\selectfont\corrigesChapitre{4}{La division : entière et décimale}
\begin{blocLivre}{cor-4-1}\begin{corrige}{1}{Une division posée}\demarche{Ne pas oublier le zéro du quotient.}
Quotient $\res{209}$, reste $\res4$ : $1\,258=6\times209+4$ et $4<6$. Contrôle : $6\times209+4=1\,258$.
\end{corrige}\end{blocLivre}
\begin{blocLivre}{cor-4-2}\begin{corrige}{2}{Le zéro au quotient}\demarche{Abaisser chiffre à chiffre sans sauter de rang.}
$3\,618:9=\res{402}$ : le 1 abaissé ne suffit pas, d’où le 0. $2\,814:7=\res{402}$ aussi : coïncidence instructive, à contrôler : $402\times7=2\,814$.
\end{corrige}\end{blocLivre}
\begin{blocLivre}{cor-4-3}\begin{corrige}{3}{Après la virgule}\demarche{Transformer le reste en rang suivant.}
$27:4$ donne 6, reste 3 ; en dixièmes : $30:4=7$, reste 2 ; en centièmes : $20:4=5$. Donc $27:4=\res{6{,}75}$. Contrôle : $6{,}75\times4=27$.
\end{corrige}\end{blocLivre}
\begin{blocLivre}{cor-4-4}\begin{corrige}{4}{Dividende décimal}\demarche{Placer la virgule au passage des dixièmes.}
$8:6=1$ reste 2 ; avec le 7 : $27:6=4$ dixièmes, reste 3 ; en centièmes : $30:6=5$. Donc $8{,}7:6=\res{1{,}45}$ ; contrôle : $1{,}45\times6=8{,}7$.
\end{corrige}\end{blocLivre}
\begin{blocLivre}{cor-4-5}\begin{corrige}{5}{Moitiés et quarts décimaux}\demarche{Diviser par 2, une ou deux fois.}
$9:2=\res{4{,}5}$ : la moitié de 9. $9:4=\res{2{,}25}$ : le quart, la moitié de la moitié.
\end{corrige}\end{blocLivre}
\begin{blocLivre}{cor-4-6}\begin{corrige}{6}{Choisir le bon quotient}\demarche{Se demander si la grandeur se coupe.}
L’argent se partage exactement : $75:4=\res{18{,}75}$ euros chacun. Les cahiers restent entiers : $75=4\times18+3$, soit $\res{18}$ cahiers par carton et 3 cahiers restants. Le contexte choisit le quotient.
\end{corrige}\end{blocLivre}
\begin{blocLivre}{cor-4-7}\begin{corrige}{7}{Le reste et l’égalité}\demarche{Comparer le reste au diviseur.}
$507=8\times62+11$ est exacte, mais $11>8$ : un paquet de plus est possible. La division euclidienne s’écrit $\res{507=8\times63+3}$, avec $3<8$.
\end{corrige}\end{blocLivre}
\begin{blocLivre}{cor-4-8}\begin{corrige}{8}{Division par dix, cent, mille}\demarche{Décaler les rangs vers la droite.}
Chaque chiffre descend de rangs : $\res{45}$ ; $\res{3{,}8}$ ; $\res{7{,}2}$.
\end{corrige}\end{blocLivre}
\begin{blocLivre}{cor-4-9}\begin{corrige}{9}{Encadrer le quotient}\demarche{Chercher deux multiples voisins du diviseur.}
$7\times35=245$ et $7\times36=252$ : $\res{35<250:7<36}$. La division donne quotient $\res{35}$ et reste $\res5$.
\end{corrige}\end{blocLivre}
\begin{blocLivre}{cor-4-10}\begin{corrige}{10}{L’erreur de virgule}\demarche{Contrôler tout quotient par la multiplication.}
$2{,}04\times6=12{,}24\ne14{,}4$ : faux. Les 24 dixièmes donnent 4 dixièmes au quotient, pas 0 puis 4 : $14{,}4:6=\res{2{,}4}$ ; contrôle : $2{,}4\times6=14{,}4$.
\end{corrige}\end{blocLivre}
\begin{blocLivre}{cor-4-11}\begin{corrige}{11}{La course en relais}\demarche{Poursuivre la division après la virgule.}
$15:4=\res{3{,}75}$ km chacun. Contrôle : $3{,}75\times4=15$ km. Le quotient décimal convient, car une distance se partage exactement.
\end{corrige}\end{blocLivre}
\begin{blocLivre}{cor-4-12}\begin{corrige}{12}{Les packs d’eau}\demarche{Faire parler le reste.}
$1\,350=8\times168+6$ : $\res{168}$ packs et $\res6$ bouteilles seules. Il manque $8-6=\res2$ bouteilles pour un pack de plus.
\end{corrige}\end{blocLivre}
\begin{blocLivre}{cor-4-13}\begin{corrige}{13}{Le ruban partagé}\demarche{Penser le dividende en dixièmes.}
$9{,}6:8=\res{1{,}2}$ m, car 96 dixièmes partagés en 8 font 12 dixièmes. En centimètres : $\res{120}$ cm.
\end{corrige}\end{blocLivre}
\begin{blocLivre}{cor-4-14}\begin{corrige}{14}{La cagnotte des quatre}\demarche{Convertir le reste en centimes.}
$58:4=\res{14{,}50}$ euros chacun : 14 euros, puis 2 euros restants qui font 200 centimes, soit 50 centimes chacun. Contrôle : $14{,}50\times4=58$ euros.
\end{corrige}\end{blocLivre}
\begin{blocLivre}{cor-4-15}\begin{corrige}{15}{Un partage impossible ?}\demarche{Arrêter la division au rang demandé et le dire.}
La division ne s’arrête pas : $10:3=3{,}33\ldots$ ; aucun nombre à deux décimales ne convient exactement. Au centimètre près, chaque morceau mesure environ $\res{3{,}33}$ m, et il reste un centimètre à répartir : $3\times3{,}33=9{,}99$. Certains partages n’ont pas d’écriture décimale exacte.
\end{corrige}\end{blocLivre}
